\section{Introduction}\label{sec:intro}

Mortality rises with age in a similar way across species, while lifespans differ by orders of magnitude
\citep{jones2014diversity,raz2026damage}. Stochastic models of damage and death connect these patterns to
mechanism \citep{strehler1960general,sacher1962stochastic,ledberg2020exponential,fedichev2025minimal}. The
saturating-removal (SR) model is one such model: damage is produced at a rate that rises with age, removed by a
process that saturates, and perturbed by noise, and death occurs when damage crosses a threshold
\citep{karin2019senescent}. It was introduced to explain senescent-cell dynamics and the Gompertz law in mice
\citep{karin2019senescent}. Since then it has been used to explain the incidence of age-related diseases
\citep{katzir2021senescent}, parabiosis experiments \citep{karin2021parabiosis}, the timing of death in
\emph{E.~coli} \citep{yang2023damage}, interventions that compress morbidity
\citep{yang2025compression}, the heritability of human lifespan \citep{shenhar2026heritability} and aging regimes
across species \citep{raz2026damage}.
% Yang et al. 2019 (Sci Adv, E. coli temporal scaling) is not cited (Naveh, 2026-11-23).

These studies set the model's parameters in different ways. Only two of them had measurements of the damage itself
\citep{karin2019senescent,yang2023damage}; the others relied on survival, incidence or mortality curves, which are
available for many species, strains and interventions. The methods that used only such curves had not been
standardized or tested thoroughly in the parameter regimes that span a wide range of organisms, or they were very
costly to run.

The cross-species study fitted the model to the survival curves of many organisms with a simulation-based
likelihood: Langevin simulation of many individuals, a kernel density estimate of the death-time distribution, and Markov-chain Monte Carlo
\citep{raz2026damage,foremanmackey2013emcee}.
The main drawback of that approach is its cost, not its answer. In an unpublished comparison on the calibrated
organisms of the cross-species study, an earlier and less accurate deterministic method, a least-squares fit of the
Fokker--Planck survival curve to the Kaplan--Meier curve, already reproduced the MCMC results: the identifiable parameter combinations agreed to a median of \C{fpmcmc-combo-dex} decades, and the
external hazard fell inside the MCMC's 95\,\% interval for \C{fpmcmc-exth-in} of \C{fpmcmc-exth-n} organisms. The fit took
about \C{fpmcmc-fp-s}\,s per survival curve on a laptop, against about \C{fpmcmc-mcmc-h}\,h on 16 cores for the MCMC.
The exact likelihood used here also ranks parameter values like the simulation-based objective (Spearman correlation
\C{simlik-spearman} over \C{simlik-finite-sim} of \C{simlik-cloud-n} points of a synthetic cohort; the objective was
minus infinity at the others), at \C{simlik-ms-fp}\,ms per evaluation against \C{simlik-ms-sim}\,ms. At the cost of the MCMC, a thorough test of the fitting method, on thousands of synthetic cohorts with known
parameters, was not feasible. The simulation-based likelihood is also noisy, and censoring, left truncation, interval-recorded deaths and an external
hazard have to be added by hand. Finally, it leaves two practical questions open: which parameters a survival curve
actually determines, and how far the reported intervals can be trusted.

Here we describe the methods now used for SR fits, and the evidence for them:
\begin{enumerate}[leftmargin=1.5em,itemsep=1pt]
  \item an exact likelihood, obtained by solving the Fokker--Planck equation of the model for the first-passage
    distribution, with an observation model for interval-recorded deaths, censoring, left truncation, a fit window,
    an external hazard and parameter changes at a given age (Section~\ref{sec:likelihood});
  \item a maximum-likelihood calibrator with profile-likelihood intervals and explicit claim shapes (a value, a bound
    or no claim), in four variants (Section~\ref{sec:estimation});
  \item what survival data identify and what they do not, and the reporting rules that follow
    (Section~\ref{sec:identifiability});
  \item a procedure for asking which parameter differs between two groups, including individuals housed together in cages, vials or plates
    (Section~\ref{sec:effects});
  \item validation: numerical accuracy against high-precision references and simulation; a benchmark of the
    calibrator on \C{v6-n-cohorts} synthetic cohorts with known parameters; and the available evidence on the error
    rates of the group comparisons (Section~\ref{sec:validation});
  \item open-source software with pinned versions, so that a study can cite the exact method it used
    (Section~\ref{sec:software}).
\end{enumerate}
% The two citing studies (mouse, mosquito) are added as examples of use in an updated version once they are public (Naveh 2026-11-22).

